


\textbf{Aqui basicamente tengo que dividirlo en varias partes. Primeramente hablar del diseño y la eleccion como 
de la arquitectura de los agentes, el por qué de cada protocolo de comunicacion, como funciona un flujo de un 
autobus llegando a la parada, el tema del suavizado de cambios entre valores de franjas, etc.
Tambien discutir sobre los dos modelos existentes, que el que incluye pasajeros puede ser considerado más 
"de agentes" por su comportamiento mas parecido a la realidad, con pasajeros que piden bajar, buses que rechazan 
subir más gente, y el otro que es más orientado a simulaciones rápidas(15 mins) pero que ambos generan mismos
resultados sobre el papel. Puedo tmb meter una ontologia? Para mostrar un poco la arquitectura de los distintos
elementos que componen cada simulacion. Hablar tambien de como se recogen los resultados, los logs etc}